%==============================================================================
% Exemplo de Proposta de Dissertação com a classe texufpel
%
% A proposta tem o mesmo tema do exemplo de dissertação
% (dissertação/exemplo-diss.tex): é o projeto que, depois de executado,
% resultaria naquele trabalho. Troque os dados do preâmbulo (os comentários
% "ALTERE" indicam o que mudar) e o texto dos capítulos.
%
% Compilação (a partir da pasta raiz do repositório):
%   pdflatex dissertação/exemplo-diss-proposta
%   bibtex   exemplo-diss-proposta
%   pdflatex dissertação/exemplo-diss-proposta
%   pdflatex dissertação/exemplo-diss-proposta
%
% As páginas citadas nos comentários são as do Manual de normas UFPel para
% trabalhos acadêmicos (2023).
%==============================================================================

% Opção diss-proposta: a proposta não é depositada na biblioteca, por isso a
% classe não gera capa (a opção capa é ignorada) nem a referência do trabalho
% acima do resumo. A antiga opção nocipinfo não é mais necessária aqui: a
% classe já a aplica às propostas.
\documentclass[diss-proposta]{texufpel}

%------------------------------------------------------------------------------
% Dados institucionais. ALTERE para o seu programa.
% Unidade feminina (Faculdade, Escola): \unidade[a]{...}.
%------------------------------------------------------------------------------
\unidade{Centro de Desenvolvimento Tecnológico}
\programa{Programa de Pós-Graduação em Computação}
\curso{Ciência da Computação}

%------------------------------------------------------------------------------
% Dados do trabalho. ALTERE.
% Os nomes abaixo são de grandes nomes da Computação e servem só de exemplo.
% Substitua pelos nomes da sua orientação.
%------------------------------------------------------------------------------
\title{Roteamento da coleta seletiva em grafos dirigidos: um estudo de caso em Pelotas}

\author{Lovelace}{Ada}                          % ALTERE: autor(a)
\advisor[Prof.ª~Dr.ª]{Liskov}{Barbara}          % ALTERE: orientador(a)
\coadvisor[Prof.~Dr.]{Dijkstra}{Edsger W.}      % ALTERE: coorientador(a), opcional

% Palavras-chave: iniciais minúsculas (exceto nomes próprios), separadas por
% ponto e vírgula, a última terminada em ponto (Manual, p. 15).
\keyword{algoritmo de Dijkstra;}
\keyword{caminho mínimo;}
\keyword{coleta seletiva;}
\keyword{grafos dirigidos.}

\begin{document}

% Orientadora: troque o rótulo da folha de rosto.
\renewcommand{\advisorname}{Orientadora}

\maketitle   % folha de rosto (sem capa, por ser proposta)
\sloppy      % evita linhas que ultrapassam a margem direita

%------------------------------------------------------------------------------
% Resumo (no máximo uma página): problema, objetivo e contribuição esperada.
%------------------------------------------------------------------------------
\begin{abstract}
A coleta seletiva de Pelotas segue rotas definidas pela experiência das
equipes, sem apoio de ferramentas de otimização. Em bairros com muitas ruas de
mão única, a ordem em que o caminhão visita os pontos de coleta pode aumentar
a distância rodada por dia. Esta proposta de dissertação trata o problema como
um problema de caminho mínimo em grafos dirigidos, em que o sentido das ruas é
representado pela direção das arestas. A malha viária de quatro setores de
coleta será modelada dessa forma, e o algoritmo de Dijkstra calculará a menor
rota entre a garagem e a cooperativa de reciclagem, passando pelos pontos de
cada setor. O objetivo principal é verificar se as rotas calculadas são mais
curtas do que as rotas praticadas hoje e medir quanto do ganho se deve à
representação correta das ruas de mão única. A contribuição esperada é um
protótipo em Python, aplicado aos quatro setores, e uma comparação, setor a
setor, entre as rotas calculadas e as atuais.
\end{abstract}

%==============================================================================
\chapter{Motivação}
% Entre uma e duas páginas.
%==============================================================================

O problema do caminho mínimo acompanha a Computação desde o fim da década de
1950. A nota de três páginas em que \citet{Dijkstra:1959:NTP} apresentou seu
algoritmo resolve o caso de origem única com pesos não negativos, e outras
soluções da mesma época, como a de Bellman, tratam pesos negativos
\citep{Cormen:2022:IA}. Hoje esses algoritmos estão em sistemas de navegação,
no roteamento de redes de computadores e em boa parte dos problemas de
logística \citep{Knuth:1997:ACP,Cormen:2022:IA}.

A coleta seletiva é um desses problemas. O caminhão parte da garagem, passa
pelos pontos de coleta e leva o material até a cooperativa de reciclagem; cada
quilômetro a mais é combustível, tempo da equipe e desgaste do veículo. Em
Pelotas, como em muitas cidades de porte médio, as rotas são definidas pela
experiência de quem dirige. Isso funciona em traçados simples, mas tende a
gerar desvios no centro histórico, onde boa parte das ruas é de mão única.

É justamente a mão única que motiva esta dissertação. Um modelo em que as ruas
são arestas sem direção calcula rotas que o caminhão não pode seguir, e a
economia que ele promete é, em parte, ilusória. Representar a malha como um
grafo dirigido corrige isso, mas também muda o resultado esperado: é possível
que a rota mínima real seja bem mais longa do que a calculada no modelo
simplificado. Medir essa diferença é a principal contribuição acadêmica do
trabalho.

O tema também se liga a outros trabalhos da Universidade Federal de Pelotas que usam a Computação para
apoiar decisões ambientais da cidade, como o de \citet{Cruz:2015:FCS}, sobre a
simulação da dispersão de poluentes na atmosfera.

%==============================================================================
\chapter{Objetivos e Resultados}
% Entre uma e três páginas.
%==============================================================================

O objetivo geral é verificar se o algoritmo de Dijkstra, aplicado a um modelo
em grafo dirigido da malha viária, reduz a distância percorrida nas rotas de
coleta seletiva de Pelotas.

Os objetivos específicos são:

\begin{enumerate}
  \item modelar a malha viária de quatro setores de coleta como grafos
        ponderados e dirigidos, com o sentido de cada rua;
  \item implementar o algoritmo de Dijkstra e validá-lo em grafos pequenos,
        com resultado conhecido;
  \item calcular as rotas mínimas de cada setor nos modelos dirigido e não
        dirigido;
  \item comparar as rotas calculadas com as rotas praticadas pela equipe de
        coleta.
\end{enumerate}

Os resultados esperados são:

\begin{enumerate}
  \item um protótipo em Python, de código aberto, que recebe o grafo de um
        setor e devolve a rota mínima;
  \item os grafos dos quatro setores, conferidos com a equipe de coleta;
  \item uma comparação, setor a setor, entre as rotas praticadas e as
        calculadas nos dois modelos;
  \item a submissão de ao menos um artigo a evento da área.
\end{enumerate}

As listas acima usam alíneas, como pede o Manual (p. 71): o texto que as
antecede termina em dois pontos, cada alínea termina em ponto e vírgula e a
última, em ponto.

%==============================================================================
\chapter{Metodologia}
% Entre uma e três páginas.
%==============================================================================

O trabalho será desenvolvido em quatro etapas, uma para cada objetivo
específico.

Na primeira, os setores serão escolhidos com a equipe de coleta, de modo a
incluir traçados diferentes: o centro, com quarteirões regulares e muitas ruas
de mão única, e bairros com loteamentos mais recentes. Cada setor será
modelado como um grafo, com a garagem, os pontos de coleta e a cooperativa como
vértices e os trechos de rua como arestas, com a distância em quilômetros como
peso. Trechos de mão dupla terão duas arestas, uma em cada sentido; trechos de
mão única, uma só.

Na segunda etapa, o algoritmo de Dijkstra será implementado em Python, com a
fila de prioridade do módulo \texttt{heapq}, seguindo a descrição de
\citet{Cormen:2022:IA}. A implementação será validada em grafos pequenos, com
distâncias mínimas calculadas à mão.

Na terceira, o protótipo calculará a rota mínima de cada setor duas vezes: no
modelo dirigido e em uma versão do mesmo grafo que ignora o sentido das ruas.
Na quarta, as rotas calculadas serão comparadas com as rotas praticadas,
informadas pela equipe de coleta, pela distância total percorrida.

%==============================================================================
\chapter{Cronograma}
% Atividades em doze meses, incluindo seminário de andamento, entrega e
% apresentação da dissertação.
%==============================================================================

As atividades previstas são:

\begin{enumerate}
  \item revisão da literatura sobre caminho mínimo e roteamento;
  \item levantamento dos setores e das rotas praticadas, com a equipe de
        coleta;
  \item modelagem dos grafos dirigidos dos setores;
  \item implementação e validação do algoritmo de Dijkstra;
  \item cálculo das rotas nos dois modelos e comparação com as rotas
        praticadas;
  \item escrita da dissertação;
  \item seminário de andamento;
  \item submissão de artigo a evento da área;
  \item entrega da dissertação à banca;
  \item apresentação da dissertação (defesa).
\end{enumerate}

A Tabela~\ref{tab:cronograma} distribui as atividades ao longo de 2026,
identificadas pela letra da alínea.

% Cronograma em tabela: uma coluna para a atividade (a letra da alínea) e uma
% para cada mês, com os meses abreviados conforme o Anexo A do Manual (p. 79).
\begin{table}[htbp]
  \caption{Cronograma de atividades em 2026}
  \label{tab:cronograma}
  \centering
  \setlength{\tabcolsep}{4pt}
  \begin{tabular}{l*{12}{c}}
    \hline
    Atividade & jan. & fev. & mar. & abr. & maio & jun. & jul. & ago. & set. & out. & nov. & dez. \\
    \hline
    a) & X & X & X &   &   &   &   &   &   &   &   &   \\
    b) &   & X & X & X &   &   &   &   &   &   &   &   \\
    c) &   &   &   & X & X &   &   &   &   &   &   &   \\
    d) &   &   &   &   & X & X & X &   &   &   &   &   \\
    e) &   &   &   &   &   &   &   & X & X &   &   &   \\
    f) &   &   & X & X & X & X & X & X & X & X & X &   \\
    g) &   &   &   &   &   &   & X &   &   &   &   &   \\
    h) &   &   &   &   &   &   &   &   &   & X &   &   \\
    i) &   &   &   &   &   &   &   &   &   &   & X &   \\
    j) &   &   &   &   &   &   &   &   &   &   &   & X \\
    \hline
  \end{tabular}
  \caption*{Fonte: Elaborada pela autora (2026).}
\end{table}

%------------------------------------------------------------------------------
% Referências: estilo abnt.bst e base bibliografia.bib, na raiz. Só as obras
% citadas no texto entram na lista.
%------------------------------------------------------------------------------
\bibliographystyle{abnt}
\bibliography{bibliografia}

%------------------------------------------------------------------------------
% Assinaturas da proponente e da orientadora.
%------------------------------------------------------------------------------
\chapter*{Assinaturas}

\vspace{2cm}

\begin{center}
\rule{8cm}{.3mm}

Ada Lovelace\\               % ALTERE
Proponente
\end{center}

\vspace{3cm}

\begin{center}
\rule{8cm}{.3mm}

Prof.ª Dr.ª Barbara Liskov\\  % ALTERE
Orientadora
\end{center}

\end{document}
